% Werke X/1 (1917), native224–226, printed217–219.
\begingroup
\makeatletter\tagsleft@true\let\veqno\@@leqno\makeatother
\clearpage\SourceLeaf{224}\NoSourceHead
\SourceGroup{IX}
\HistoricTitle{[DIE BEZIEHUNG ZWISCHEN DEN UNENDLICH VIELEN WERTEN DES ARITHMETISCH-GEOMETRISCHEN MITTELS. UNMITTELBARE ANWENDUNG DER THEORIE AUF DIE ELLIPTISCHEN TRANSZENDENTEN UND AUF DIE REKTIFIKATION DER ELLIPSE.]}
\NumberedItem{1}
\HistoricCitation{[Aus Handbuch 19, Be, Kleine Aufsätze aus verschiedenen Theilen der Mathematik, Angefangen im May 1809, S.\ 6.]}
\par\indent Medium arithm[etico-geometricum] inter \(A\) et \(B=M\)
\[\begin{array}{ll}
a=A,&b=\surd(AA-BB)=C\\
a'=\frac12(a+b),&b'=\surd ab\\
a''=\frac12(a'+b'),&b''=\surd a'b'
\end{array}\]
\begin{center}.\quad .\quad .\quad .\quad .\quad .\quad .\quad .\quad .\end{center}
\begin{center}Medium [inter \(a\) et \(b\)] \(=m\).\end{center}
{\small
\begin{equation*}\tag*{[1]}
M=\frac{k\pi m}{\log16+2\log\frac{a'}B-\log\frac{a'}{a''}-\frac12\log\frac{a''}{a'''}-\frac14\log\frac{a'''}{a^{\mathrm{IV}}}-\ldots}
\end{equation*}}
\begin{equation*}\tag*{[2]}
=\frac{\frac12k\pi m}{\log4+\log\frac{a'}B-\frac12\log\frac{a'}{a''}-\text{etc.}}
\end{equation*}
\begin{equation*}\tag*{[3]}\log\tfrac12k\pi=9{,}8339042\end{equation*}
\begin{equation*}\tag*{[4]}
\left.\begin{aligned}
c'c'+2c''c''+4c'''c'''+\text{etc.}\\
+C'C'+2C''C''+4C'''C'''+\text{etc.}
\end{aligned}\right\}=\frac12aa-\frac{2mM}{\pi}.
\end{equation*}
\ClosingRule
\par\bigskip\noindent\small X1.\hfill28\normalsize
\clearpage\SourceLeaf{225}\ArithmeticHead{218}{L}{ANALYSIS.\quad NACHLASS.}
\HistoricCitation{[Ein Blatt in Fg, Kapsel 46a, Wasserzeichen FHF 1810.]}
\HistoricTitle{Arithmetisch-geometrische Mittel.}
\NumberedItem{2}
\[\begin{aligned}
m&=\mu(1+2x+2x^4+2x^9+\text{etc.})^2\\
n&=\mu(1-2x+2x^4-2x^9+\text{etc.})^2\\
\surd(mm-nn)&=\mu(2x^{\frac14}+2x^{\frac94}+2x^{\frac{25}4}+\text{etc.})^2
\end{aligned}\]
\par\indent Das AGMittel zwischen \(m\) und \(\surd(mm-nn)\) sei \(=\lambda\). Man hat dann
\[\lambda=\frac{\pi\mu}{\log\frac1x}\quad\text{oder}\quad x=e^{-\frac{\pi\mu}{\lambda}}.\]
\noindent Für \textsc{Brigg}ische Logarithmen wird \(\lambda=\dfrac{a\mu}{\log\frac1x}\), wo \(\log a=0{,}1349342\).
\par\indent Auch kann man sich folgender Formel bedienen:
\[\begin{aligned}
x&=\frac{mm-nn}{16m'.\surd m''.{}^4\!\surd m'''.{}^8\!\surd m^{\mathrm{IV}}.\mathbin{\&}\mathrm c}\\[.5em]
&=\frac{m-n}{8\surd m''.{}^4\!\surd m'''.{}^8\!\surd m^{\mathrm{IV}}\ldots}\mathpunct{\raisebox{.5ex}{.}}
\end{aligned}\]
\par\smallskip\hrule\vspace{2pt}\hrule
\begin{center}\textls[120]{Beispiel}\end{center}
{\small\setlength{\tabcolsep}{10pt}\renewcommand{\arraystretch}{1.18}
\begin{center}\begin{tabular}{rrrr}
3,0000000&0,4771213&3,0000000&0,4771213\\
1,0000000&0,0000000&2,8284270&0,4515450\\\hline
2,0000000&0,3010300&2,9142135&0,4645214\\
1,7320508&0,2385606&2,9129510&0,4643332\\\hline
1,8660254&0,2709175&2,9135822&0,4644273\\
1,8612098&0,2697953&&\\\hline
1,8636176&0,2703568&0,2703567&9,6989700\\
1,8636159&0,2703564&0,2709175&0,5716671\\[-.6em]
\multicolumn{2}{c}{\rule{12.5em}{.4pt}}&\rule{5.5em}{.4pt}&\rule{5.5em}{.4pt}\\[-.5em]
1,8636167&0,2703566&0,2706371&9,1273029\\\cline{1-2}\noalign{\vskip2pt}\cline{1-2}
&&0,3010300&
\end{tabular}\end{center}}
\clearpage\SourceLeaf{226}\ArithmeticHead{219}{R}{\scriptsize DIE UNENDLICH VIELEN WERTE DES ARITHM.-GEOM.\ MITTELS.}
\noindent\begin{minipage}[t]{.48\linewidth}
\[\begin{array}{lr}
a\ .\ .\ .\ .&0{,}1349342\\
\mu\ .\ .\ .\ .&0{,}2703566\\
\mathrm C.\lambda\ .\ .\ .&9{,}5355727\\\hline
&9{,}9408635\\
&0{,}8726970\\
x\ .\ .\ .\ .&9{,}1273030
\end{array}\]
\end{minipage}\hfill\begin{minipage}[t]{.48\linewidth}
\[0{,}6396307=\text{Canon }\left[=\frac\mu\lambda\right]\]
\end{minipage}
\par\medskip\hrule
\NumberedItem{3}
\par\indent Die AGMittel gestalten sich anders, wenn man für ein \(n'\;n''\;n'''\;\&\mathrm c\) den negativen Werth wählt: doch sind alle Resultate in folgender Form begriffen:
\begin{equation*}\tag*{[1]}
\frac1{(\mu)}=\frac1\mu+\frac{4ik}\lambda\mathpunct{\raisebox{.5ex}{.}}
\end{equation*}
\begin{center}\textls[120]{Beispiel für einen imaginären Werth des AGMittels.}\end{center}
{\small\setlength{\arraycolsep}{4pt}\renewcommand{\arraystretch}{1.25}
\[\begin{array}{rrr rrr}
3{,}0000000&&0{,}4771213&0^0&&\\
1{,}0000000&&0{,}0000000&360^0&&\\\hline
2{,}0000000&&0{,}3010300&0^0&&\\
-1{,}7320508&&0{,}2385606&180^0&&\\\hline
0{,}1339746&&9{,}1270225&0^0&&\\
&+1{,}8612098i&0{,}2697953&90^0&&\\\hline
0{,}0669873&+0{,}9306049i&9{,}9698876&85^0&52'&58''{,}10\\
0{,}3530969&+0{,}3530969i&9{,}6984089&45&0&0\\\hline
\overline0{,}2100421&+0{,}6418509i&9{,}8295254&71&52&46{,}58\\
0{,}2836930&+0{,}6208239i&9{,}8341482&65&26&29{,}05\\\hline
0{,}2468676&+0{,}6313374i&9{,}8311572&68&38&36{,}05\\
0{,}2470649&+0{,}6324002i&9{,}8318368&68&39&37{,}82\\\hline
0{,}2469962\frac12&+0{,}6318688i&9{,}8314971&68&39&6{,}95\\
0{,}2469962\frac12&+0{,}6318685i&9{,}8314970&68&39&6{,}93\\\hline
0{,}2469962\frac12&+0{,}6318686\frac12i&9{,}8314970\frac12&68&39&6{,}94
\end{array}\]}
\[\frac1{(\mu)}=+0{,}5365910-1{,}3728774i=\frac1\mu+\frac{4i}\lambda.\]
\par\bigskip\hfill\small28\textsuperscript{*}\normalsize
\clearpage\endgroup
